	\section{InsertionSort e MergeSort}

\begin{softnote}
	Ora studiamo due algoritmi di ordinamento fondamentali. InsertionSort è semplice ma lento, MergeSort è più complesso ma veloce. Capirai la differenza tra approccio \textbf{incrementale} e \textbf{divide et impera}.
\end{softnote}

\subsubsection{InsertionSort}

Inserisce ogni elemento nella posizione corretta tra quelli già ordinati
(come le carte da gioco).

\begin{nota}
	\textbf{Analogia con le carte:} quando ordini una mano di carte, prendi una carta alla volta e la inserisci nella posizione giusta tra quelle già ordinate. Questo è esattamente quello che fa InsertionSort!
\end{nota}

\begin{algorithm}
	\caption{InsertionSort}
	\KwIn{$A = \langle a_1, \ldots, a_n \rangle$, array di elementi 
		confrontabili con $\leq$}
	\KwOut{$A$ ordinato rispetto a $\leq$}
	\For{$i = 2$ \KwTo $n$}{
		$key = A[i]$ \tcp*{Inserisce $A[i]$ nella parte ordinata $a_1, \ldots, a_{i-1}$}
		$j = i - 1$\;
		\While{$j > 0$ \textbf{and} $A[j] > key$}{
			$A[j+1] = A[j]$\;
			$j = j - 1$\;
		}
		$A[j+1] = key$\;
	}
	\Return{$A$}
\end{algorithm}

\begin{softnote}
	\textbf{Come funziona InsertionSort passo p\\asso:}
	
	\begin{enumerate}
		\item Parte dall'elemento in posizione 2 (il primo è già ``ordinato'')
		\item Prende l'elemento corrente (\texttt{key})
		\item Lo confronta con gli elementi precedenti (da destra a sinistra)
		\item Sposta gli elementi più grandi verso destra
		\item Inserisce \texttt{key} nella posizione corretta
		\item Ripete per tutti gli elementi
	\end{enumerate}
\end{softnote}

\begin{nota}
	Ordina gli elementi \textbf{sul posto} (in-place), cioè non usa array aggiuntivi ma modifica direttamente l'array di input.
\end{nota}

\textbf{Ciclo Algoritmo}

\begin{center}
	
	\vspace{0.4em}
	\begin{tabular}{|c|c|c|c|c|}
		\hline
		5 & 2 & 4 & 6 & 1 \\
		\hline
	\end{tabular}
	$= A$
	
	\vspace{0.4em}
	{\small\textit{Ogni blocco rappresenta un ciclo for (0 o più cicli while)}}
	\vspace{0.4em}
	
	\scalebox{0.9}{%
		\begin{tabular}{|c|c|c|c|c|l|}
			\hline
			\textit{1} & \textit{2} & \textit{3} & \textit{4} & \textit{5} & \textit{Commento} \\
			\hline
			5 & \textcolor{codered}{2} & 4 & 6 & 1 &
			\multirow{2}{*}{\parbox{3cm}{\small Seq. ord.: 5 \\ Scambio 5$\to$2}} \\
			j & i & & & & \\[3pt]
			\hline
			2 & 5 & \textcolor{codered}{4} & 6 & 1 &
			\multirow{2}{*}{\parbox{3cm}{\small Seq. ord.: 2,5 \\ Scambio 5$\to$4}} \\
			$\leftarrow$ & j & i & & & \\[3pt]
			\hline
			2 & 4 & 5 & \textcolor{codered}{6} & 1 &
			\multirow{2}{*}{\parbox{3cm}{\small Seq. ord.: 2,4,5 \\ Nessun scambio}} \\
			& $\leftarrow$ & j & i & & \\[3pt]
			\hline
			2 & 4 & 5 & 6 & \textcolor{codered}{1} &
			\multirow{2}{*}{\parbox{3cm}{\small Seq. ord.: 2,4,5,6 \\ 1 $\to$ pos. I}} \\
			& & $\leftarrow$ & j & i & \\[3pt]
			\hline
			1 & 2 & 4 & 5 & 6 &
			\multirow{2}{*}{\parbox{3cm}{\small Seq. ordinata}} \\
			& & & & i & \\[3pt]
			\hline
		\end{tabular}%
	}
\end{center}

\begin{softnote}
	Nella tabella:
	\begin{itemize}
		\item \textcolor{codered}{Rosso} = elemento che stiamo inserendo (\texttt{key})
		\item $i$ = indice elemento corrente
		\item $j$ = indice di confronto (parte da $i-1$ e va a ritroso)
		\item Le frecce $\leftarrow$ mostrano il movimento degli elementi verso destra
	\end{itemize}
\end{softnote}

\pagebreak
\subsubsection*{Invariante di ciclo}

\begin{definizione}[Invariante di ciclo]
	Un'asserzione (proprietà booleana) sempre vera all'inizio
	di ogni iterazione di un ciclo.
\end{definizione}

\begin{nota}
	L'invariante di ciclo è come una \textbf{promessa} che l'algoritmo mantiene: ``prima di ogni giro del ciclo, questa proprietà è vera''. Serve per dimostrare matematicamente che l'algoritmo funziona correttamente.
\end{nota}

La validità si dimostra in tre passi:
\begin{enumerate}
	\item \textbf{Inizializzazione:} è vera prima della prima iterazione.
	\item \textbf{Conservazione:} se è vera prima di un'iterazione,
	rimane vera prima di quella successiva.
	\item \textbf{Conclusione:} quando il ciclo termina, l'invariante
	fornisce una proprietà utile a dimostrare la correttezza.
\end{enumerate}

\begin{esempio}[Invariante di InsertionSort]
	\textit{``All'inizio del ciclo \texttt{for}, il sottoarray
		$A[1:i-1]$ è ordinato e formato dagli stessi elementi
		che erano in $A[1:i-1]$.''}
	
	\begin{itemize}
		\item \textbf{Inizializzazione:} $i=2$, quindi $A[1:1]$
		è un solo elemento $\Rightarrow$ ordinato. \checkmark
		\item \textbf{Conservazione:} \texttt{key}$= A[i]$,
		il ciclo \texttt{while} sposta $A[i-1], A[i-2],\ldots$
		a destra finché non trova la posizione corretta per
		\texttt{key}. Quindi $A[1:i]$ diventa ordinato con
		gli stessi elementi. All'iterazione successiva
		$i$ aumenta di 1, l'invariante vale su $A[1:i-1]$.
		\checkmark
		\item \textbf{Conclusione:} il ciclo termina con $i = n+1$.
		Per l'invariante, $A[1:n]$ è ordinato con gli stessi
		elementi $\Rightarrow$ l'algoritmo è corretto. \checkmark
	\end{itemize}
\end{esempio}

\begin{softnote}
	\textbf{In parole semplici:}
	\begin{itemize}
		\item \textbf{Inizializzazione}: all'inizio funziona (un elemento è sempre ordinato)
		\item \textbf{Conservazione}: se funziona al passo $i$, funziona anche al passo $i+1$
		\item \textbf{Conclusione}: alla fine, l'intero array è ordinato
	\end{itemize}
	
	È come l'effetto domino: se il primo pezzo cade e ogni pezzo fa cadere il successivo, alla fine cadono tutti!
\end{softnote}

\subsubsection*{Modello RAM}


Per analizzare la complessità si usa un modello astratto di
computer: la \textbf{Random Access Machine (RAM)}.

\begin{itemize}
	\item Un solo processore con memoria ad accesso casuale.
	\item Numeri interi e a virgola mobile con dimensione fissa
	(es.\ 64 bit per parola).
	\item Istruzioni elementari (lettura, scrittura, somma,
	sottrazione, moltiplicazione, divisione, confronto,
	\texttt{if}, salto, chiamata a procedura).
	\item Ogni istruzione ha \textbf{costo costante} $= 1$.
\end{itemize}

\begin{nota}
	È una semplificazione: nella realtà cache, pipeline e
	parallelismo influenzano i tempi reali. In questo corso
	si ignorano questi dettagli.
	
	\textbf{Perché usiamo il modello RAM?} Perché ci permette di confrontare algoritmi in modo uniforme, indipendentemente dall'hardware specifico su cui girano.
\end{nota}

\subsubsection*{Dimensione dell'input e tempo di esecuzione}

\begin{definizione}[Dimensione dell'input]
	Misura della quantità di bit/word necessari per rappresentare
	l'input. Un indicatore è \textbf{corretto} se la dimensione
	reale è in relazione polinomiale con esso.
\end{definizione}

\begin{definizione}[Tempo di esecuzione]
	Numero totale di operazioni elementari eseguite dall'algoritmo
	nel modello RAM, in funzione della dimensione $n$ dell'input.
\end{definizione}

\begin{nota}
	\textbf{Indicatore corretto:} una misura semplificata della dimensione dell'input che mantiene una relazione ragionevole con la dimensione reale.
	
	\textbf{Esempi:}
	\begin{itemize}
		\item Per un array di $n$ numeri: indicatore = $n$ (numero di elementi)
		\item Per una matrice $n \times n$: indicatore = $n$ (anche se contiene $n^2$ elementi)
		\item Per un numero intero: indicatore = numero di bit necessari a rappresentarlo
	\end{itemize}
\end{nota}

Esempi di indicatori corretti:
\begin{itemize}
	\item \textbf{Array di $n$ elementi:} indicatore $= n$. \checkmark
	\item \textbf{Matrice $n \times n$:} indicatore $= n$
	(la dimensione reale $n \cdot k(n+1)$ è polinomiale in $n$).
	\checkmark
	\item \textbf{Intero $n$:} indicatore $= \lfloor \lg n \rfloor + 1$
	(numero di bit). \textbf{Non} si può usare $n$ come
	indicatore perché $2^{|I|} \approx n$, relazione
	esponenziale. \checkmark
\end{itemize}

\begin{esempio}[Matrice $n \times n$]
	Sia $A$ una matrice $n \times n$ di interi, dove ogni intero
	richiede $k$ bit. La RAM memorizza la matrice riga per riga.
	
	Ogni riga occupa $nk$ bit per gli interi, più $k$ bit per il
	marcatore di fine riga. L'intera matrice richiede:
	\[
	|I| = n\bigl(n+1\bigr)k \text{ bit}
	= n(kn + k) \text{ bit}
	= \boxed{kn^2 + kn \text{ bit}}
	= \Theta(n)
	\]
	
	\begin{center}
		\begin{tikzpicture}[scale=1.2]
			% --- matrice n x n ---
			\draw[step=0.6cm, gray!60, thin] (0,0) grid (3.6, 3.6);
			\draw[thick] (0,0) rectangle (3.6, 3.6);
			\node[above] at (1.8, 3.6) {$n$};
			\node[left]  at (0,  1.8)  {$n$};
			\node[font=\tiny] at (0.3, 3.3) {$a_{1,1}$};
			\node[font=\tiny] at (3.3, 3.3) {$a_{1,n}$};
			\node[font=\tiny] at (0.3, 0.3) {$a_{n,1}$};
			\node[font=\tiny] at (3.3, 0.3) {$a_{n,n}$};
			
			% --- freccia RAM ---
			\draw[->, thick] (4.0, 1.8) -- (5.2, 1.8);
			\node[above, font=\small] at (4.6, 1.9) {RAM};
			
			% --- stack RAM ---
			\foreach \y in {0, 0.4, 0.8, 1.2, 1.6, 2.0, 2.4, 2.8, 3.2}
			\draw[draw=gray!80, fill=codebg]
			(5.4, \y) rectangle (7.2, \y+0.35);
			\node[font=\tiny] at (6.3, 3.35) {$a_{1,1}$};
			\node[font=\tiny] at (6.3, 2.95) {$a_{1,2}$};
			\node[font=\tiny] at (6.3, 2.55) {$\vdots$};
			\node[font=\tiny] at (6.3, 2.15) {$a_{1,n}$};
			\node[font=\tiny, color=codered] at (6.3, 1.75) {\texttt{fine riga} ($k$ bit)};
			\node[font=\tiny] at (6.3, 1.35) {$a_{2,1}$};
			\node[font=\tiny] at (6.3, 0.95) {$\vdots$};
			\node[font=\tiny] at (6.3, 0.55) {$a_{n,n}$};
			\node[font=\tiny, color=codered] at (6.3, 0.15) {\texttt{fine riga} ($k$ bit)};
			
			% --- graffa dimensione ---
			\draw[decorate, decoration={brace, amplitude=5pt, mirror}]
			(7.3, 0) -- (7.3, 3.55)
			node[midway, right=6pt, font=\small] {$n$ righe};
			
			% --- annotazione dimensione riga ---
			\draw[decorate, decoration={brace, amplitude=4pt}]
			(5.4, 3.6) -- (7.2, 3.6)
			node[midway, above=5pt, font=\tiny]
			{$nk + k$ bit per riga};
		\end{tikzpicture}
	\end{center}
	
	Si usa come indicatore $|I| = n$. È \textbf{corretto} perché
	la dimensione reale $kn^2 + kn$ è un \textbf{polinomio in $n$}.
\end{esempio}

\begin{nota}
	L'errore comune è usare il \emph{valore} di un intero $n$ come
	dimensione dell'input invece del numero di bit necessari a
	rappresentarlo. Questo porta a sottostimare la complessità.
	
	\textbf{Esempio:} il numero 1000 ha valore 1000, ma occupa solo $\lceil \log_2(1000) \rceil = 10$ bit. La dimensione dell'input è 10, non 1000!
\end{nota}

\subsubsection*{Analisi di InsertionSort}

$c_k$ = costo esecuzione istruzione alla riga $k$. \quad
$t_i - 1$ = \# cicli \textbf{while} nel $i$-esimo ciclo \textbf{for}.

\begin{center}
	\begin{tabular}{rll r r}
		\toprule
		& \textsc{Insertion-Sort}$(A, n)$ & & \textit{cost} & \textit{times} \\
		\midrule
		1 & \textbf{for} $i = 2$ \textbf{to} $n$ & & $c_1$ & $n$ \\
		2 & \quad $key = A[i]$ & & $c_2$ & $n-1$ \\
		3 & \quad \textcolor{codegreen}{\texttt{// Insert $A[i]$ into the sorted subarray $A[1:i-1]$}} & & $0$ & $n-1$ \\
		4 & \quad $j = i - 1$ & & $c_4$ & $n-1$ \\
		5 & \quad \textbf{while} $j > 0$ \textbf{and} $A[j] > key$ & & $c_5$ & $\displaystyle\sum_{i=2}^{n} t_i$ \\[6pt]
		6 & \qquad $A[j+1] = A[j]$ & & $c_6$ & $\displaystyle\sum_{i=2}^{n} (t_i - 1)$ \\[6pt]
		7 & \qquad $j = j - 1$ & & $c_7$ & $\displaystyle\sum_{i=2}^{n} (t_i - 1)$ \\[6pt]
		8 & \quad $A[j+1] = key$ & & $c_8$ & $n-1$ \\
		\bottomrule
	\end{tabular}
\end{center}

\begin{softnote}
	\textbf{Contare le iterazioni:}
	
	\textbf{for} $i = \ell$ \textbf{to} $n$ \quad $\Rightarrow$ \quad $\#\text{Cicli} = n - \ell + 1$
	
	\textbf{Esempio:} \textbf{for} $i = 2$ \textbf{to} $n$ \quad $\Rightarrow$ \quad $n - 1$ iterazioni
\end{softnote}


\[
T(n) = c_1 n + c_2(n-1) + c_4(n-1)
+ c_5 \sum_{i=2}^{n} t_i
+ c_6 \sum_{i=2}^{n} (t_i-1)
+ c_7 \sum_{i=2}^{n} (t_i-1)
+ c_8(n-1)
\]


Indichiamo con $c_k$ il costo dell'istruzione alla riga $k$,
e con $t_i$ il numero di volte che il ciclo \texttt{while} viene
eseguito durante l'$i$-esimo ciclo \texttt{for}.

\begin{nota}
	$t_i$ dipende dai dati:
	\begin{itemize}
		\item Se $A[i]$ è già al posto giusto: $t_i = 1$ (un solo confronto)
		\item Se $A[i]$ deve andare all'inizio: $t_i = i$ (confronta con tutti)
	\end{itemize}
\end{nota}

\pagebreak

\subsubsection*{Caso migliore: array già ordinato}

Se l'array è già ordinato, il corpo del \texttt{while} non viene
mai eseguito: $t_i = 1$ per ogni $i$.
Le somme $\sum(t_i - 1) = 0$, quindi:

\[
T(n) = c_1 n + c_2(n-1) + c_4(n-1) + c_5(n-1) + c_8(n-1)
\]
\[
= \underbrace{(c_1 + c_2 + c_4 + c_5 + c_8)}_{a} n
- \underbrace{(c_2 + c_4 + c_5 + c_8)}_{b}
\]
\[
= an - b \quad \Rightarrow \quad \boxed{T(n) = \Theta(n)}
\]

\begin{nota}
	Nel caso migliore la crescita è \textbf{lineare}: ogni elemento
	viene confrontato una sola volta con il suo predecessore e
	nessuno scambio viene effettuato.
	
	\textbf{Esempio pratico:} se l'array è $[1, 2, 3, 4, 5]$, l'algoritmo fa solo $n-1$ confronti e nessuno spostamento. Velocissimo!
\end{nota}

\subsubsection*{Caso peggiore: array ordinato al contrario}

Se l'array è in ordine decrescente, ogni $A[i]$ deve essere
confrontato con \textbf{tutti} gli elementi in $A[1:i-1]$,
quindi $t_i = i$. Usando $\sum_{i=2}^{n} i = \frac{n(n+1)}{2} - 1$:

\[
T(n) = c_1 n + c_2(n-1) + c_4(n-1)
+ c_5\!\left(\frac{n(n+1)}{2} - 1\right)
+ (c_6 + c_7)\frac{n(n-1)}{2}
+ c_8(n-1)
\]

Raccogliendo i termini per grado:

\[
T(n) = \underbrace{\frac{c_5 + c_6 + c_7}{2}}_{a} n^2
+ \underbrace{\left(c_1 + c_2 + c_4 + \frac{c_5 - c_6 - c_7}{2} + c_8\right)}_{b} n
- \underbrace{(c_2 + c_4 + c_5 + c_8)}_{c}
\]
\[
\Rightarrow \quad \boxed{T(n) = \Theta(n^2)}
\]

\begin{nota}
	Nel caso peggiore la crescita è \textbf{quadratica}: ogni nuovo
	elemento deve ``scendere'' fino in fondo alla parte già ordinata,
	confrontandosi con tutti gli elementi precedenti.
	
	\textbf{Esempio pratico:} se l'array è $[5, 4, 3, 2, 1]$:
	\begin{itemize}
		\item Il 4 fa 1 confronto
		\item Il 3 fa 2 confronti
		\item Il 2 fa 3 confronti
		\item L'1 fa 4 confronti
	\end{itemize}
	Totale: $1 + 2 + 3 + 4 = 10 = \frac{n(n-1)}{2}$ confronti. Molto lento!
\end{nota}

\begin{center}
	\begin{tabular}{lll}
		\toprule
		Caso & Input & Complessità \\
		\midrule
		Migliore & Array già ordinato      & $\Theta(n)$   \\
		Peggiore & Array ordinato inverso  & $\Theta(n^2)$ \\
		\bottomrule
	\end{tabular}
\end{center}

\begin{hardnote}
	Solitamente si considera come complessità di un algoritmo il
	tempo di esecuzione con input che rappresenta il caso peggiore.
	
	\begin{enumerate}[label=\arabic*.]
		\item È una garanzia sulla descrizione del comportamento dell'algoritmo.
		\item In alcuni casi, il caso peggiore è il più frequente.
		\item Il caso medio spesso è simile al caso peggiore (vedi InsertionSort).
	\end{enumerate}
	
	\textbf{Per InsertionSort:} il caso medio è $\Theta(n^2)$ come il caso peggiore, perché in media metà degli elementi vanno spostati.
\end{hardnote}

\pagebreak

\subsection{MergeSort}

MergeSort è un algoritmo di ordinamento basato sul metodo \textbf{divide et impera},
più efficiente di InsertionSort nel caso peggiore: $\Theta(n \lg n)$ contro $\Theta(n^2)$.

\subsubsection{Metodo Divide et Impera}

Il metodo divide et impera è un paradigma di progettazione degli algoritmi
che si articola in tre fasi ricorsive:

\begin{itemize}
	\item \textbf{Divide:} il problema di dimensione $n$ viene suddiviso
	in $a$ sotto-problemi di dimensione $n/b$.
	\item \textbf{Impera:} ogni sotto-problema viene risolto ricorsivamente.
	La ricorsione si ferma quando il problema è sufficientemente piccolo
	da essere risolto direttamente (caso base).
	\item \textbf{Combina:} le soluzioni dei sotto-problemi vengono
	unite per ottenere la soluzione del problema originale.
\end{itemize}

\begin{nota}
	InsertionSort usa un approccio \textbf{incrementale}: estende la
	soluzione di $A[1 \ldots i-1]$ inserendo $A[i]$ nella posizione corretta.
	MergeSort usa invece \textbf{divide et impera}: risolve due metà
	indipendentemente e le fonde.
\end{nota}

\subsubsection{Descrizione dell'algoritmo}

Applicando divide et impera al problema dell'ordinamento:

\begin{itemize}
	\item \textbf{Divide:} l'array $A[p \ldots r]$ viene spezzato in due
	sotto-array $A[p \ldots q]$ e $A[q+1 \ldots r]$, dove
	$q = \lfloor (p+r)/2 \rfloor$.
	\item \textbf{Impera:} i due sotto-array vengono ordinati ricorsivamente
	con MergeSort. Quando $p = r$ (sotto-array di un solo elemento),
	la ricorsione termina.
	\item \textbf{Combina:} i due sotto-array ordinati vengono fusi
	dall'operazione \textsc{Merge}$(A, p, q, r)$.
\end{itemize}

L'operazione chiave è \textsc{Merge}: dati due sotto-array adiacenti già
ordinati, li fonde in un unico sotto-array ordinato.

\subsubsection{Pseudocodice}

Dato $A[1 \ldots n]$, MergeSort si invoca con
\textsc{Merge-Sort}$(A,\, 1,\, n)$.

\begin{algorithm}[H]
	\caption{\textsc{Merge-Sort}$(A, p, r)$}
	\KwIn{Array $A$, indici $p \leq r$}
	\KwOut{$A[p \ldots r]$ ordinato in loco}
	\BlankLine
	\If{$p < r$}{
		$q = \lfloor (p + r) / 2 \rfloor$\;
		\textsc{Merge-Sort}$(A,\, p,\, q)$\;
		\textsc{Merge-Sort}$(A,\, q+1,\, r)$\;
		\textsc{Merge}$(A,\, p,\, q,\, r)$\;
	}
\end{algorithm}

\begin{algorithm}[H]
	\caption{\textsc{Merge}$(A, p, q, r)$}
	\KwIn{$A[p \ldots q]$ e $A[q+1 \ldots r]$ già ordinati}
	\KwOut{$A[p \ldots r]$ ordinato}
	\BlankLine
	$n_L = q - p + 1$\;
	$n_R = r - q$\;
	crea array $L[0 \ldots n_L-1]$ e $R[0 \ldots n_R-1]$\;
	\BlankLine
	\For{$i = 0$ \KwTo $n_L - 1$}{
		$L[i] = A[p + i]$\;
	}
	\For{$j = 0$ \KwTo $n_R - 1$}{
		$R[j] = A[q + 1 + j]$\;
	}
	\BlankLine
	$i = 0$;\quad $j = 0$;\quad $k = p$\;
	\BlankLine
	\While{$i < n_L$ \normalfont\textbf{and} $j < n_R$}{
		\eIf{$L[i] \leq R[j]$}{
			$A[k] = L[i]$;\; $i = i + 1$\;
		}{
			$A[k] = R[j]$;\; $j = j + 1$\;
		}
		$k = k + 1$\;
	}
	\BlankLine
	\While{$i < n_L$}{
		$A[k] = L[i]$;\; $i = i + 1$;\; $k = k + 1$\;
	}
	\While{$j < n_R$}{
		$A[k] = R[j]$;\; $j = j + 1$;\; $k = k + 1$\;
	}
\end{algorithm}

\begin{nota}
	Questa è la versione più semplice di \textsc{Merge}. Nel Cormen (ed.\ IV)
	e in rete esistono versioni più compatte che usano la tecnica della
	\textbf{sentinella}: si aggiunge $+\infty$ in fondo a $L$ e $R$,
	eliminando i due cicli \textbf{while} finali.
\end{nota}


\subsubsection{Correttezza: Invariante di ciclo di Merge}

\begin{definizione}[Invariante di \textsc{Merge}]
	All'inizio di ogni iterazione del ciclo \textbf{while} principale,
	$A[p \ldots k-1]$ contiene, ordinati, gli $i+j$ elementi minori
	di tutti gli elementi in $L[0 \ldots n_L-1]$ e $R[0 \ldots n_R-1]$.
	$L[i]$ e $R[j]$ sono gli elementi più piccoli dei rispettivi array
	non ancora copiati in $A$.
\end{definizione}

\begin{proof}
	\begin{enumerate}
		\item \textbf{Inizializzazione:} $k = p$, quindi $A[p \ldots k-1]$
		è vuoto ($0$ elementi copiati). Con $i = j = 0$, $L[0]$ e
		$R[0]$ sono i minimi non ancora copiati. L'invariante è vera. \checkmark
		
		\item \textbf{Conservazione:} supponiamo $L[i] \leq R[j]$ (il caso
		$L[i] > R[j]$ è simmetrico). $L[i]$ è il minimo tra tutti gli
		elementi non ancora copiati, quindi viene copiato in $A[k]$.
		Incrementando $k$ e $i$, l'array $A[p \ldots k-1]$ contiene
		ora $i+j$ elementi ordinati e $L[i]$ è il nuovo minimo
		non copiato di $L$. L'invariante è mantenuta. \checkmark
		
		\item \textbf{Conclusione:} il ciclo termina quando $i = n_L$
		oppure $j = n_R$, cioè quando uno dei due array è esaurito.
		I rimanenti elementi dell'altro array (già ordinati e tutti
		maggiori di quelli copiati) vengono aggiunti dai due cicli
		\textbf{while} finali. Al termine, $A[p \ldots r]$ è ordinato
		e contiene gli stessi elementi originali. \checkmark
	\end{enumerate}
\end{proof}

\subsubsection{Complessità}

\paragraph{Complessità di \textsc{Merge}.}

Sia $n = r - p + 1$ la dimensione del sotto-array da fondere:
\begin{itemize}
	\item Copiare in $L$ e $R$: $n_L + n_R = n$ operazioni
	\item Ciclo \textbf{while} principale: al più $n$ iterazioni
	\item Cicli \textbf{while} finali: al più $n$ iterazioni totali
\end{itemize}
\[
T_{\textsc{Merge}}(n) = \Theta(n)
\]

\paragraph{Equazione di ricorrenza di MergeSort.}

Con $a = 2$ sotto-problemi di dimensione $n/2$, costo di divisione
$D(n) = \Theta(1)$ e costo di fusione $C(n) = \Theta(n)$:
\[
T(n) = \begin{cases}
	\Theta(1)                & \text{se } n = 1           \\[4pt]
	2\,T(n/2) + \Theta(n)   & \text{altrimenti}
\end{cases}
\]

Applicando il \textbf{Teorema dell'Esperto}:
\[
\boxed{T(n) = \Theta(n \lg n)} \qquad \text{dove } \lg = \log_2
\]

\paragraph{Soluzione per albero di ricorsione.}

\begin{center}
	\begin{nota}
		Ad ogni livello $i$ dell'albero vi sono $2^i$ sotto-problemi di
		dimensione $n/2^i$, ciascuno con costo di fusione $\Theta(n/2^i)$.
		Il costo totale per livello è:
		\[
		2^i \cdot \Theta\!\left(\frac{n}{2^i}\right) = \Theta(n)
		\]
		Il numero di livelli è $\lg n + 1$ (le foglie hanno dimensione 1,
		con $n/2^{\lg n} = 1$). Sommando su tutti i livelli:
		\[
		T(n) = (\lg n + 1) \cdot \Theta(n) = \Theta(n \lg n)
		\]
	\end{nota}
\end{center}

\subsubsection{Confronto con InsertionSort}

\begin{nota}
	I due algoritmi rappresentano due filosofie opposte:
	InsertionSort usa un approccio \textbf{incrementale}, MergeSort usa
	\textbf{divide et impera}. La scelta dipende dal contesto.
\end{nota}

\begin{center}
	\begin{tabular}{lll}
		\toprule
		\textbf{Proprietà}      & \textbf{InsertionSort}    & \textbf{MergeSort}        \\
		\midrule
		Paradigma               & Incrementale              & Divide et impera          \\
		Caso migliore           & $\Theta(n)$               & $\Theta(n \lg n)$         \\
		Caso peggiore           & $\Theta(n^2)$             & $\Theta(n \lg n)$         \\
		Caso medio              & $\Theta(n^2)$             & $\Theta(n \lg n)$         \\
		Memoria extra           & $\Theta(1)$ (in-place)    & $\Theta(n)$ (array $L$, $R$) \\
		Stabile                 & Sì                        & Sì                        \\
		Adatto a input piccoli  & Sì                        & No                        \\
		\bottomrule
	\end{tabular}
\end{center}

\begin{softnote}
	\textbf{Quando usare InsertionSort?}
	\begin{itemize}
		\item Array piccoli ($n \leq 20\text{--}30$): la costante nascosta
		in $\Theta(n^2)$ è talmente piccola che batte MergeSort in pratica.
		\item Array \textbf{quasi ordinati}: con pochi elementi fuori posto,
		InsertionSort è quasi $\Theta(n)$.
		\item Ambienti con memoria limitata: nessun array aggiuntivo.
	\end{itemize}
	
	\textbf{Quando usare MergeSort?}
	\begin{itemize}
		\item Array grandi: $\Theta(n \lg n)$ diventa rapidamente più
		conveniente di $\Theta(n^2)$.
		\item Quando si vuole una \textbf{garanzia} sul caso peggiore
		(non dipende dall'ordine dell'input).
		\item Ordinamento esterno (dati su disco): MergeSort si adatta
		naturalmente alla fusione di blocchi.
	\end{itemize}
\end{softnote}

\begin{hardnote}
	\textbf{Perché $\Theta(n \lg n)$ è asintoticamente ottimale?}
	
	Qualsiasi algoritmo di ordinamento \textbf{basato su confronti}
	richiede nel caso peggiore almeno $\Omega(n \lg n)$ confronti
	(dimostrazione tramite alberi di decisione, cap.\ 8 CLRS).
	MergeSort è quindi \textbf{asintoticamente ottimale} per questa
	classe di algoritmi.
	
	InsertionSort, con il suo $\Theta(n^2)$, è lontano dall'ottimo
	per input grandi, ma rimane utile grazie alle sue costanti
	moltiplicative molto piccole e alla semplicità implementativa.
\end{hardnote}

\begin{esempio}[Crescita a confronto]
	Per $n = 10^6$ elementi:
	\begin{itemize}
		\item InsertionSort (caso peggiore): $\sim 10^{12}$ operazioni
		\item MergeSort: $\sim 2 \times 10^7$ operazioni
		($n \lg_2 n \approx 2 \times 10^7$)
	\end{itemize}
	La differenza è di \textbf{cinque ordini di grandezza}.
	Su un processore da $10^9$ operazioni/secondo, InsertionSort
	impiegherebbe $\sim 16$ minuti, MergeSort $\sim 0{,}02$ secondi.
\end{esempio}